\documentclass[10pt]{article}
\usepackage[margin=20mm]{geometry}
\usepackage{amsmath,amssymb,fontspec,unicode-math,xurl,hyperref}
\IfFontExistsTF{DejaVu Serif}{\setmainfont{DejaVu Serif}}{\setmainfont{Latin Modern Roman}}
\setmathfont{latinmodern-math.otf}
\setlength{\parindent}{0pt}\setlength{\parskip}{6pt}\sloppy
\title{Position: Strong Consumer Protection is an Inalienable Defense for AI Safety in the United States}
\author{ELVIA independent paper verification}
\date{Review date: 23 October 2026 (America/Chicago)}
\begin{document}
\maketitle
\section*{Identity, exact version and outcome}
I25-40112 / display0013. Serena Booth, Position: Strong Consumer Protection is an Inalienable Defense for AI Safety in the United States. Reviewed the official ICML 2025 proceedings PDF, PMLR 267:81118-81129; OpenReview j5eEF77JQz. Locations use printed pages 1-12, equal to one-based PDF pages. Source SHA256: 14388ff872856a8c186c0f6d0639e357babe83e182620d69af77dcc247924b53.\par
Supported outcome: Minor, limited to an overbroad attribution of a model-architecture restriction to CFPB guidance. The paper’s next-page qualification supplies substantial counterevidence and preserves its main lesson. No demonstrated Major defect or central-claim damage was established. The RealPage counterfactual and the fraud statistic are recorded separately as scope/evidence clarifications.\par
\section*{Scope, assumptions and checked argument}
Read the full 12-page paper, including the main argument on pp. 1-9 and the bibliography on pp. 9-12. This policy position paper supplies case studies rather than a quantitative causal estimate or a new algorithm. Selected contemporaneous primary sources were checked for the most material interpretation questions. No target-paper code, experimental supplement or separate appendix was identified in the inspected publisher record. No code execution, statistical reproduction, exhaustive legal audit or present-day compliance assessment was performed.\par
The broad AI-safety definition and exclusions are explicit in sections 1.1-1.3, p. 2. Gradual harm emergence is expressly speculative in section 2, pp. 2-3, and the following paragraph preserves a case for consumer protection even under abrupt emergence. Section4.1 treats AI-specific rules as potentially complementary; section 4.2 acknowledges compliance costs and regulatory capture; section 5 recognizes enforcement and political fragility. These passages prevent reading the title as a theorem that current law is sufficient or automatically enforced.\par
\section*{Problem1. Explanation requirements are overstated as a deep-learning restriction}
Source claim and exact location: section 3.2, printed p.4, final paragraph, attributes to the CFPB an effective prohibition on deep-learning systems for credit decisions, especially without significant human oversight. The suspect step replaces a functional requirement for explanations with a restriction on model architecture. Evidence E1-E3.\par
Complete check: the cited19 September 2023 CFPB announcement requires accurate, specific adverse-action reasons that reflect actual decision factors, including when complex models are used. The contemporaneous Circular 2022-03, Response and Analysis, states the prohibition conditionally: a creditor may not use a complex algorithm when it cannot supply the required specific, accurate reasons. Endnote1 discusses post-hoc explanations and requires their accuracy to be validated, recognizing that this may be infeasible for some models. Thus the relevant condition is explanatory compliance, not membership in the deep-learning model class. Neither inspected source makes human oversight a substitute for that requirement.\par
Alternative reading and counterevidence: the author may mean that many practically inscrutable deep models cannot meet the requirement. The same section on p. 5 explicitly narrows the statement to systems that remain inscrutable and says institutions must be able to provide accurate explanations. That is much closer to the official source. This audit therefore identifies a bounded local overstatement, not an invalidation of the case study or a finding about every deployed credit model.\par
Impact: the p. 4 formulation can exaggerate the claimed innovation-versus-consumer-protection tradeoff by suggesting a restriction on an entire architecture family. The central example that technology does not remove adverse-action explanation duties remains supported. Judgment: Minor. If interpreted solely as a prediction about practical model capability, the architecture-wide prediction instead Needs Evidence.\par
Precise source-supported clarification: creditors using complex algorithms remain subject to adverse-action explanation requirements, and may not use a system when doing so prevents them from providing the specific, accurate reasons required for those actions. This restores the conditional scope and agrees with the manuscript’s later qualification. This is a verified source-based clarification, not a proof that any particular credit system complies with all applicable law. The manuscript was not edited and author adoption has not been requested.\par
\section*{Separate clarifications and unresolved claims}
\subsection*{Supporting checks}
Section3.1, pp. 3-4, proposes that replacing private inputs with publicly inferred signals could preserve rent-setting performance and make legal applicability more ambiguous. The original DOJ complaint, filed23 August 2024, alleges information-sharing agreements in paragraphs 225-234 and broader agreements to align pricing in paragraphs 236-241. The latter are not defined only by the privacy of inputs. Complaint printedpp.81-85 correspond to one-based PDF pages86-90. These are allegations, not adjudicated findings.\par
A public-data substitution alone does not show that a common pricing arrangement lacks an agreement. However, the paper says the case becomes more ambiguous, not that public data guarantees legality. A fair evaluation distinguishes independent responding algorithms from an agreed common pricing service. The suggestion that inferred data would maintain comparable performance also lacks a measured inference error or downstream sensitivity analysis. Both are Needs Evidence within a speculative counterfactual; neither establishes a false reported experiment. No finding about RealPage’s ultimate liability is made.\par
\subsection*{Fraud: the measured increase is monetary loss}
Section3.5, p. 6, refers to a 30\% increase in fraud/scams and nearly9 billion dollars of losses. The cited FTC release reports more than30\% growth in reported dollar losses, reaching nearly8.8 billion dollars in 2022. The official explanatory post separately says fraud reports decreased from2.9 million to 2.4 million. The rounded dollar amount in the paper is reasonable; the wording should identify reported monetary loss as the statistic. It does not establish a 30\% rise in incident counts, population prevalence, or AI-attributable fraud. Because the surrounding sentence discusses losses, this is retained as metric clarification rather than a separate confirmed defect.\par
\subsection*{Remaining case-study coverage}
Sections3.3-3.6 on disparate impact, supervisory authority, transfers and privacy were read for argumentative dependence and acknowledged gaps. Their full statutory exceptions, enforcement histories and all economic estimates were not independently audited. No quantitative effectiveness estimate was validated. Later administrative developments are not used retroactively as evidence of an original2025 error.\par
\section*{Overall dependence, version search and conclusion}
The core thesis remains a conditional case for preserving, enforcing and improving technology-neutral protections. It does not require the local architecture-wide interpretation of CFPB guidance. The inspected PMLR record, exact-title searches and an author publication listing did not identify a separate later manuscript, code release, supplement or formal erratum. OpenReview modification metadata is not evidence that a particular substantive correction occurred. No corrected historical issue has been relabeled as a current defect.\par
Item | Judgment | Repair verification | Adoption / remaining work\par
CFPB description | Minor local overstatement | Conditional source-based clarification supported | Not adopted; no original manuscript changed\par
RealPage legal/performance counterfactual | Needs Evidence / scope clarification | No legal or empirical repair asserted | Separate agreement structure; evaluate public-data performance\par
30\% fraud statistic | Metric clarification | Use reported dollar losses | No inference about population incidence or AI attribution\par
Central consumer-protection thesis | No central damage demonstrated | No overall repair required by these checks | Scientific/legal human judgment pending\par
\section*{Consulted references}
E1. ICML 2025PMLR proceedings;12 pages: https://proceedings.mlr.press/v267/booth25a.html\par
E2. CFPB 19 September 2023 announcement: https://www.consumerfinance.gov/archive/newsroom/cfpb-issues-guidance-on-credit-denials-by-lenders-using-artificial-intelligence/\par
E3. CFPB Circular 2022-03: https://www.consumerfinance.gov/compliance/circulars/circular-2022-03-adverse-action-notification-requirements-in-connection-with-credit-decisions-based-on-complex-algorithms/\par
E4. Original complaint filed23 August 2024;Case1:24-cv-00710Document1: https://www.justice.gov/opa/media/1364976/dl?inline\par
E5. FTC 23 February 2023 release: https://www.ftc.gov/news-events/news/press-releases/2023/02/new-ftc-data-show-consumers-reported-losing-nearly-88-billion-scams-2022\par
E6. FTC 2023 official explanatory post: https://www.ftc.gov/business-guidance/blog/2023/02/ftc-crunches-2022-numbers-see-where-scammers-continue-crunch-consumers\par
\section*{Supporting checks}
> Within assigned staging root:
> python.exe sources/fetch\_paper.py 0013 [publisher-linked PDF URL]
> pdftotext.exe -layout papers/0013/0013.pdf sources/0013.txt
> python.exe sources/acquire\_0013\_evidence.py
> Source records: sources/0013-acquisition.json; sources/0013-evidence/acquisition.json.
> Original PDF:12 pages,199457 bytes; SHA256 stated above.
> DOJ complaint snapshot:4509811 bytes; SHA2567dd6882be591d55d4767dee6385eb7fab824736a37dff43b647f3b7baa89a7a1.\par
The CFPB 2023 announcement and full 2022 circular were successfully read through official web pages, but local snapshot requests returned HTTP403. The FTC release was successfully read on the web; its local download returned HTTP404. These failures are recorded rather than treated as successful snapshots. A delegated attempt to open the archived full 2023 circular reached an Archive-It rate limit; this report does not claim to have read that circular in full. The downloaded DOJ complaint and the source-paper PDF succeeded.\par
No human review or author approval is asserted. Coverage, findings and provenance are separate per-paper files.\par
\end{document}
